

Subbab ini merinci spesifikasi setiap \textit{audit event} yang telah
diidentifikasi pada Tabel~\ref{tab:threat-event-mapping}. Untuk setiap
\textit{event}, dijelaskan konteks aktivitas yang dicatat, kemungkinan nilai
pada atribut umum, serta atribut tambahan beserta kemungkinan nilainya.

Atribut umum yang terdapat pada seluruh jenis \textit{audit event} adalah
\texttt{source\_component}, \texttt{actor\_id}, \texttt{target\_object},
\texttt{outcome}, dan \texttt{action\_type}. Kemungkinan nilai dari atribut
\texttt{outcome} dan \texttt{action\_type} berlaku seragam di seluruh
\textit{event} dan dijelaskan pada Tabel~\ref{tab:nilai-outcome} dan
Tabel~\ref{tab:nilai-action-type}.

\begin{table}[htbp]
\centering
\caption{Kemungkinan Nilai Atribut \texttt{outcome}}
\label{tab:nilai-outcome}
\renewcommand{\arraystretch}{1.3}
\begin{tabular}{|p{2.5cm}|p{9.5cm}|}
\hline
\textbf{Nilai} & \textbf{Kondisi} \\
\hline
\texttt{Success} &
Aktivitas berhasil diselesaikan sesuai dengan yang dimaksudkan. \\
\hline
\texttt{Failure} &
Aktivitas gagal diselesaikan akibat kesalahan teknis, kondisi yang tidak
terpenuhi, atau kegagalan pada komponen yang dihubungi (mis. IOTA, IPFS,
Redis tidak merespons). \\
\hline
\texttt{Denied} &
Aktivitas ditolak secara eksplisit karena tidak memenuhi persyaratan
otorisasi, misalnya \textit{capability} tidak valid, JWT kadaluarsa, atau
peran tidak sesuai. \\
\hline
\end{tabular}
\end{table}

\begin{table}[htbp]
\centering
\caption{Kemungkinan Nilai Atribut \texttt{action\_type}}
\label{tab:nilai-action-type}
\renewcommand{\arraystretch}{1.3}
\begin{tabular}{|p{2.5cm}|p{9.5cm}|}
\hline
\textbf{Nilai} & \textbf{Kondisi} \\
\hline
\texttt{CREATE} &
Membuat atau menginisialisasi objek, entri, atau sumber daya baru.
Contoh: mendaftarkan akun, menerbitkan \textit{activation key}, membuat
entri rekam medis, mengunggah konten ke IPFS. \\
\hline
\texttt{READ} &
Membaca atau mengambil data dari objek yang sudah ada tanpa melakukan
perubahan. Contoh: membaca rekam medis, mengambil data administratif,
mengunduh konten dari IPFS. \\
\hline
\texttt{UPDATE} &
Memperbarui atau memodifikasi objek yang sudah ada. Contoh: memperbarui
entri rekam medis, memperbarui \textit{activation key}, memperbarui
metadata administratif. \\
\hline
\texttt{DELETE} &
Menghapus atau mencabut objek, entri, atau izin yang sudah ada. Contoh:
mencabut delegasi akses (\texttt{revoke\_access}), \textit{sign out},
menghapus \textit{nonce} dari Redis setelah digunakan. \\
\hline
\texttt{VALIDATE} &
Memverifikasi keabsahan, keaslian, atau kelengkapan suatu objek atau
klaim tanpa mengubah \textit{state}. Contoh: memvalidasi QR \textit{payload},
memvalidasi JWT, memvalidasi \textit{capability on-chain}. \\
\hline
\texttt{EXECUTE} &
Menjalankan proses yang telah dipersiapkan sebelumnya. Berbeda dari
\texttt{CREATE} karena objeknya sudah ada dan tinggal dieksekusi. Contoh:
mengeksekusi transaksi yang telah disponsori \textit{gas}. \\
\hline
\end{tabular}
\end{table}

% -----------------------------------------------------------------------
\subsubsection{EV1 --- Authentication}
\label{subsec:ev1-authentication}

EV1 mencatat setiap percobaan autentikasi pengguna pada \textit{Patient Client}
maupun \textit{Hospital Client}, baik yang berhasil maupun yang gagal. Aktivitas
yang tercakup meliputi proses \textit{sign-in} menggunakan \textit{activation key},
\textit{patient ID}, atau \textit{seed words}, serta proses \textit{sign-out}.
Pencatatan \textit{event} ini diperlukan untuk mendeteksi pola percobaan autentikasi
yang mencurigakan, seperti percobaan berulang yang mengindikasikan serangan
\textit{brute force}.

\begin{table}[htbp]
\centering
\caption{Atribut Umum EV1 Authentication}
\label{tab:ev1-umum}
\renewcommand{\arraystretch}{1.3}
\begin{tabular}{|p{3.5cm}|p{8.5cm}|}
\hline
\textbf{Atribut} & \textbf{Kemungkinan Nilai} \\
\hline
\texttt{source\_component} &
\texttt{HospitalTauriBackend} | \texttt{PatientTauriBackend} \\
\hline
\texttt{actor\_id} &
IOTA \textit{address} pengguna yang melakukan autentikasi. \\
\hline
\texttt{target\_object} &
\texttt{activation\_key} | \texttt{patient\_id} | \texttt{seed\_words}
--- menyatakan mekanisme autentikasi yang digunakan. \\
\hline
\texttt{action\_type} &
\texttt{VALIDATE} | \texttt{CREATE} | \texttt{DELETE} --- \texttt{VALIDATE}
untuk verifikasi kredensial, \texttt{CREATE} untuk inisialisasi sesi,
\texttt{DELETE} untuk \textit{sign-out}. \\
\hline
\texttt{outcome} &
\texttt{Success} | \texttt{Failure} | \texttt{Denied} \\
\hline
\end{tabular}
\end{table}

\begin{table}[htbp]
\centering
\caption{Atribut Tambahan EV1 Authentication}
\label{tab:ev1-detail}
\renewcommand{\arraystretch}{1.3}
\begin{tabular}{|p{4cm}|p{7.5cm}|}
\hline
\textbf{Atribut} & \textbf{Kemungkinan Nilai dan Keterangan} \\
\hline
\texttt{auth\_method} &
\texttt{activation\_key} | \texttt{patient\_id} | \texttt{seed\_words}
--- metode spesifik yang digunakan dalam percobaan autentikasi. \\
\hline
\texttt{role} &
\texttt{Pasien} | \texttt{PersonnelMedis} | \texttt{PersonnelAdministratif}
| \texttt{AdminFasyankes} --- peran pengguna yang sedang melakukan autentikasi. \\
\hline
\texttt{attempt\_count} &
Bilangan bulat positif --- jumlah percobaan autentikasi yang telah dilakukan
pada sesi atau dalam periode tertentu. \\
\hline
\texttt{failure\_reason} &
\texttt{invalid\_credentials} | \texttt{account\_not\_found} |
\texttt{session\_expired} | \texttt{null} apabila \texttt{outcome} bukan
\texttt{Failure}. \\
\hline
\texttt{device\_info} &
Informasi perangkat yang digunakan, misalnya \textit{fingerprint} perangkat
atau platform sistem operasi. \\
\hline
\end{tabular}
\end{table}

% -----------------------------------------------------------------------
\subsubsection{EV2 --- QR Validation}
\label{subsec:ev2-qr-validation}

EV2 mencatat proses validasi \textit{payload} QR Access Delegation yang dipindai
oleh pasien dari QR \textit{code} yang ditampilkan oleh \textit{hospital personnel}.
\textit{Payload} QR berisi IOTA \textit{address} dan PRE \textit{public key}
\textit{personnel} dalam format \texttt{<iota\_address>@<hex\_pre\_pub\_key>}.
Pencatatan \textit{event} ini diperlukan untuk mendeteksi pemalsuan
\textit{payload} QR maupun kebocoran informasi delegasi akses.

\begin{table}[htbp]
\centering
\caption{Atribut Umum EV2 QR Validation}
\label{tab:ev2-umum}
\renewcommand{\arraystretch}{1.3}
\begin{tabular}{|p{3.5cm}|p{8.5cm}|}
\hline
\textbf{Atribut} & \textbf{Kemungkinan Nilai} \\
\hline
\texttt{source\_component} &
\texttt{PatientTauriBackend} \\
\hline
\texttt{actor\_id} &
IOTA \textit{address} pasien yang memindai QR. \\
\hline
\texttt{target\_object} &
IOTA \textit{address} \& PRE \textit{public key hospital personnel} yang
terdapat dalam QR. \\
\hline
\texttt{action\_type} &
\texttt{READ} | \texttt{VALIDATE} \\
\hline
\texttt{outcome} &
\texttt{Success} | \texttt{Failure} \\
\hline
\end{tabular}
\end{table}

\begin{table}[htbp]
\centering
\caption{Atribut Tambahan EV2 QR Validation}
\label{tab:ev2-detail}
\renewcommand{\arraystretch}{1.3}
\begin{tabular}{|p{4cm}|p{7.5cm}|}
\hline
\textbf{Atribut} & \textbf{Kemungkinan Nilai dan Keterangan} \\
\hline
\texttt{qr\_content\_hash} &
Nilai \textit{hash} SHA-256 dari \textit{payload} QR yang dipindai ---
digunakan untuk membuktikan konten yang diterima tanpa menyimpan nilai
aslinya. \\
\hline
\texttt{hospital\_personnel\_iota\_address} &
IOTA \textit{address} \textit{hospital personnel} yang terkandung dalam
QR. \\
\hline
\texttt{hospital\_personnel\_pre\_public\_key\_fingerprint} &
\textit{Fingerprint} (hash) dari PRE \textit{public key personnel} yang
terkandung dalam QR. \\
\hline
\texttt{hospital\_name} &
Nama fasyankes tempat \textit{personnel} bernaung, apabila tersedia dalam
\textit{payload}. \\
\hline
\texttt{hospital\_personnel\_name} &
Nama \textit{personnel} yang ditampilkan dalam QR, apabila tersedia. \\
\hline
\texttt{validation\_step} &
\texttt{parse} | \texttt{format\_check} | \texttt{address\_lookup} ---
tahapan validasi terakhir yang berhasil atau gagal. \\
\hline
\end{tabular}
\end{table}

% -----------------------------------------------------------------------
\subsubsection{EV3 --- Capability Creation}
\label{subsec:ev3-capability-creation}

EV3 mencatat proses pembuatan dan pencabutan \textit{capability} akses yang
diberikan pasien kepada \textit{hospital personnel}. \textit{Capability} ini
tersimpan sebagai entri \textit{on-chain} (HospitalPersonnelAccess) dan \textit{kfrag}
di Redis. Pencatatan \textit{event} ini diperlukan untuk membuktikan bahwa
pasien benar-benar telah memberikan atau mencabut delegasi akses kepada
\textit{personnel} tertentu, sehingga menangkal ancaman penyangkalan (T6).

\begin{table}[htbp]
\centering
\caption{Atribut Umum EV3 Capability Creation}
\label{tab:ev3-umum}
\renewcommand{\arraystretch}{1.3}
\begin{tabular}{|p{3.5cm}|p{8.5cm}|}
\hline
\textbf{Atribut} & \textbf{Kemungkinan Nilai} \\
\hline
\texttt{source\_component} &
\texttt{PatientTauriBackend} | \texttt{ProxyReencryption} \\
\hline
\texttt{actor\_id} &
IOTA \textit{address} pasien yang mendelegasikan akses. \\
\hline
\texttt{target\_object} &
IOTA \textit{address} \textit{hospital personnel} penerima akses. \\
\hline
\texttt{action\_type} &
\texttt{CREATE} | \texttt{DELETE} --- \texttt{CREATE} untuk pemberian
akses, \texttt{DELETE} untuk pencabutan akses. \\
\hline
\texttt{outcome} &
\texttt{Success} | \texttt{Failure} \\
\hline
\end{tabular}
\end{table}

\begin{table}[htbp]
\centering
\caption{Atribut Tambahan EV3 Capability Creation}
\label{tab:ev3-detail}
\renewcommand{\arraystretch}{1.3}
\begin{tabular}{|p{4cm}|p{7.5cm}|}
\hline
\textbf{Atribut} & \textbf{Kemungkinan Nilai dan Keterangan} \\
\hline
\texttt{access\_type} &
\texttt{read} | \texttt{update} --- cakupan akses yang diberikan. \\
\hline
\texttt{exp\_duration\_ms} &
Durasi kedaluwarsa \textit{capability} dalam milidetik. \\
\hline
\texttt{k\_frag\_fingerprint} &
\textit{Fingerprint} (hash) dari \textit{kfrag} yang disimpan di Redis,
tanpa mengekspos nilai \textit{kfrag} asli. \\
\hline
\texttt{transaction\_digest} &
\textit{Digest} transaksi IOTA yang merekam pembuatan atau pencabutan
\textit{capability on-chain}. \\
\hline
\texttt{hospital\_name} &
Nama fasyankes tempat \textit{personnel} penerima akses bernaung. \\
\hline
\texttt{nonce\_used} &
\textit{Nonce} yang digunakan dalam proses pemberian akses untuk mencegah
\textit{replay attack}. \\
\hline
\end{tabular}
\end{table}

% -----------------------------------------------------------------------
\subsubsection{EV4 --- Medical Record Access}
\label{subsec:ev4-medical-record-access}

EV4 mencatat setiap aktivitas pembacaan maupun penulisan data rekam medis
yang dieksekusi oleh PRE Server atas permintaan \textit{hospital personnel}.
Aktivitas ini mencakup pengambilan data administratif, pengambilan data klinis
beserta proses \textit{re-encryption}, pembuatan entri rekam medis baru, dan
pembaruan entri rekam medis. Pencatatan \textit{event} ini diperlukan untuk
membuktikan bahwa akses terhadap data sensitif pasien dilakukan dengan otorisasi
yang sah (T7, T13).

\begin{table}[htbp]
\centering
\caption{Atribut Umum EV4 Medical Record Access}
\label{tab:ev4-umum}
\renewcommand{\arraystretch}{1.3}
\begin{tabular}{|p{3.5cm}|p{8.5cm}|}
\hline
\textbf{Atribut} & \textbf{Kemungkinan Nilai} \\
\hline
\texttt{source\_component} &
\texttt{ProxyReencryption} \\
\hline
\texttt{actor\_id} &
IOTA \textit{address} \textit{hospital personnel} yang mengakses. \\
\hline
\texttt{target\_object} &
IPFS CID rekam medis (untuk data klinis) | IOTA \textit{address} pasien
(untuk data administratif). \\
\hline
\texttt{action\_type} &
\texttt{CREATE} | \texttt{READ} | \texttt{UPDATE} \\
\hline
\texttt{outcome} &
\texttt{Success} | \texttt{Failure} | \texttt{Denied} \\
\hline
\end{tabular}
\end{table}

\begin{table}[htbp]
\centering
\caption{Atribut Tambahan EV4 Medical Record Access}
\label{tab:ev4-detail}
\renewcommand{\arraystretch}{1.3}
\begin{tabular}{|p{4cm}|p{7.5cm}|}
\hline
\textbf{Atribut} & \textbf{Kemungkinan Nilai dan Keterangan} \\
\hline
\texttt{patient\_iota\_address} &
IOTA \textit{address} pasien pemilik data yang diakses. \\
\hline
\texttt{record\_index} &
Indeks rekam medis yang diakses atau dibuat, apabila relevan. \\
\hline
\texttt{role\_used} &
\texttt{MedicalPersonnel} | \texttt{AdministrativePersonnel} --- peran
\textit{personnel} yang tercantum pada JWT saat akses dilakukan. \\
\hline
\texttt{purpose\_used} &
\texttt{read} | \texttt{update} --- tujuan akses yang tercantum pada JWT. \\
\hline
\texttt{jwt\_sub} &
\textit{Subject} JWT yang digunakan --- biasanya IOTA \textit{address
personnel}. \\
\hline
\texttt{capability\_valid} &
\texttt{true} | \texttt{false} --- hasil validasi \textit{capability
on-chain} sebelum akses diproses. \\
\hline
\texttt{ipfs\_cid} &
CID IPFS konten klinis yang diakses atau diunggah, \texttt{null} apabila
operasi hanya menyangkut data administratif. \\
\hline
\texttt{reencryption\_performed} &
\texttt{true} | \texttt{false} --- menyatakan apakah proses
\textit{re-encryption} PRE dilakukan pada operasi ini. \\
\hline
\end{tabular}
\end{table}

% -----------------------------------------------------------------------
\subsubsection{EV5 --- Personnel Activation Key}
\label{subsec:ev5-personnel-activation-key}

EV5 mencatat aktivitas penerbitan dan pembaruan \textit{activation key} yang
dilakukan oleh Admin Fasyankes untuk \textit{personnel} di bawahnya.
\textit{Activation key} ini memungkinkan \textit{personnel} mendaftarkan diri
ke dalam sistem. Pencatatan \textit{event} ini diperlukan untuk membuktikan
bahwa admin fasyankes benar-benar menerbitkan \textit{activation key} kepada
\textit{personnel} tertentu (T8).

\begin{table}[htbp]
\centering
\caption{Atribut Umum EV5 Personnel Activation Key}
\label{tab:ev5-umum}
\renewcommand{\arraystretch}{1.3}
\begin{tabular}{|p{3.5cm}|p{8.5cm}|}
\hline
\textbf{Atribut} & \textbf{Kemungkinan Nilai} \\
\hline
\texttt{source\_component} &
\texttt{HospitalTauriBackend} \\
\hline
\texttt{actor\_id} &
IOTA \textit{address} Admin Fasyankes yang menerbitkan \textit{key}. \\
\hline
\texttt{target\_object} &
\textit{Hash} ID personel penerima \textit{activation key}. \\
\hline
\texttt{action\_type} &
\texttt{CREATE} | \texttt{UPDATE} \\
\hline
\texttt{outcome} &
\texttt{Success} | \texttt{Failure} \\
\hline
\end{tabular}
\end{table}

\begin{table}[htbp]
\centering
\caption{Atribut Tambahan EV5 Personnel Activation Key}
\label{tab:ev5-detail}
\renewcommand{\arraystretch}{1.3}
\begin{tabular}{|p{4cm}|p{7.5cm}|}
\hline
\textbf{Atribut} & \textbf{Kemungkinan Nilai dan Keterangan} \\
\hline
\texttt{admin\_iota\_address} &
IOTA \textit{address} Admin Fasyankes yang menerbitkan \textit{key}. \\
\hline
\texttt{personnel\_id\_hash} &
\textit{Hash} dari ID personel penerima \textit{activation key}. \\
\hline
\texttt{role\_assigned} &
\texttt{MedicalPersonnel} | \texttt{AdministrativePersonnel} | \texttt{Admin}
--- peran yang ditetapkan pada \textit{activation key}. \\
\hline
\texttt{hospital\_id\_hash} &
\textit{Hash} dari ID fasyankes tempat \textit{personnel} bernaung. \\
\hline
\texttt{transaction\_digest} &
\textit{Digest} transaksi IOTA yang merekam penerbitan atau pembaruan
\textit{activation key on-chain}. \\
\hline
\texttt{activation\_key\_hash} &
\textit{Hash} dari nilai \textit{activation key} yang diterbitkan. \\
\hline
\end{tabular}
\end{table}

% -----------------------------------------------------------------------
\subsubsection{EV6 --- Facility Registration}
\label{subsec:ev6-facility-registration}

EV6 mencatat aktivitas registrasi fasyankes baru beserta Admin Fasyankes-nya
yang dilakukan oleh Ministry Administrator. Pencatatan \textit{event} ini
diperlukan untuk membuktikan bahwa Kementerian benar-benar telah meregistrasi
fasyankes tertentu beserta administratornya (T9).

\begin{table}[htbp]
\centering
\caption{Atribut Umum EV6 Facility Registration}
\label{tab:ev6-umum}
\renewcommand{\arraystretch}{1.3}
\begin{tabular}{|p{3.5cm}|p{8.5cm}|}
\hline
\textbf{Atribut} & \textbf{Kemungkinan Nilai} \\
\hline
\texttt{source\_component} &
\texttt{MinistryTauriBackend} \\
\hline
\texttt{actor\_id} &
IOTA \textit{address} Ministry Administrator. \\
\hline
\texttt{target\_object} &
\textit{Hash} ID fasyankes yang diregistrasi. \\
\hline
\texttt{action\_type} &
\texttt{CREATE} | \texttt{UPDATE} \\
\hline
\texttt{outcome} &
\texttt{Success} | \texttt{Failure} \\
\hline
\end{tabular}
\end{table}

\begin{table}[htbp]
\centering
\caption{Atribut Tambahan EV6 Facility Registration}
\label{tab:ev6-detail}
\renewcommand{\arraystretch}{1.3}
\begin{tabular}{|p{4cm}|p{7.5cm}|}
\hline
\textbf{Atribut} & \textbf{Kemungkinan Nilai dan Keterangan} \\
\hline
\texttt{facility\_id} &
Identitas fasyankes yang diregistrasi. \\
\hline
\texttt{facility\_name} &
Nama fasyankes yang diregistrasi. \\
\hline
\texttt{administrator\_id} &
Identitas Admin Fasyankes yang ditetapkan saat registrasi. \\
\hline
\texttt{transaction\_digest} &
\textit{Digest} transaksi IOTA yang merekam registrasi fasyankes. \\
\hline
\texttt{activation\_key\_hash} &
\textit{Hash} dari \textit{facility activation key} yang diterbitkan
saat registrasi. \\
\hline
\end{tabular}
\end{table}

% -----------------------------------------------------------------------
\subsubsection{EV7 --- PRE Service Operation}
\label{subsec:ev7-pre-service-operation}

EV7 mencatat setiap permintaan layanan yang diterima oleh PRE Server dari
komponen \textit{client} maupun dari PRE Server itu sendiri. Aktivitas yang
tercakup meliputi seluruh \textit{endpoint} REST PRE Server: \texttt{/nonce},
\texttt{/keys}, \texttt{/medical-record}, dan \texttt{/administrative}.
Pencatatan \textit{event} ini diperlukan untuk mendeteksi penyamaran sebagai
PRE Server (\textit{man-in-the-middle}), penyadapan komunikasi, dan pembanjiran
\textit{endpoint} (T10, T11, T12).

\begin{table}[htbp]
\centering
\caption{Atribut Umum EV7 PRE Service Operation}
\label{tab:ev7-umum}
\renewcommand{\arraystretch}{1.3}
\begin{tabular}{|p{3.5cm}|p{8.5cm}|}
\hline
\textbf{Atribut} & \textbf{Kemungkinan Nilai} \\
\hline
\texttt{source\_component} &
\texttt{ProxyReencryption} \\
\hline
\texttt{actor\_id} &
IOTA \textit{address} pemohon | \texttt{proxy\_iota\_address} apabila
PRE Server bertindak sebagai aktor. \\
\hline
\texttt{target\_object} &
\textit{Endpoint} PRE Server yang dipanggil, misalnya
\texttt{/api/v1/nonce}, \texttt{/api/v1/keys},
\texttt{/api/v1/medical-record}. \\
\hline
\texttt{action\_type} &
\texttt{CREATE} | \texttt{READ} | \texttt{VALIDATE} \\
\hline
\texttt{outcome} &
\texttt{Success} | \texttt{Failure} | \texttt{Denied} \\
\hline
\end{tabular}
\end{table}

\begin{table}[htbp]
\centering
\caption{Atribut Tambahan EV7 PRE Service Operation}
\label{tab:ev7-detail}
\renewcommand{\arraystretch}{1.3}
\begin{tabular}{|p{4cm}|p{7.5cm}|}
\hline
\textbf{Atribut} & \textbf{Kemungkinan Nilai dan Keterangan} \\
\hline
\texttt{endpoint\_called} &
\textit{Path} lengkap \textit{endpoint} yang dipanggil, misalnya
\texttt{/api/v1/medical-record}. \\
\hline
\texttt{request\_id} &
Pengenal unik permintaan yang dapat digunakan untuk menghubungkan
\textit{event} ini dengan \textit{event} lain pada siklus yang sama. \\
\hline
\texttt{caller\_component} &
\texttt{PatientTauriBackend} | \texttt{HospitalTauriBackend} |
\texttt{ProxyReencryption} --- komponen yang mengajukan permintaan. \\
\hline
\texttt{channel\_encryption} &
\texttt{TLS} | \texttt{none} --- mekanisme pengamanan kanal komunikasi
yang digunakan. \\
\hline
\texttt{jwt\_purpose} &
\texttt{read} | \texttt{update} | \texttt{null} --- tujuan yang tercantum
pada JWT yang disertakan dalam permintaan, apabila ada. \\
\hline
\texttt{http\_status\_code} &
Kode status HTTP yang dikembalikan PRE Server sebagai respons, misalnya
\texttt{200}, \texttt{401}, \texttt{403}, \texttt{500}. \\
\hline
\texttt{latency\_ms} &
Durasi pemrosesan permintaan oleh PRE Server dalam milidetik. \\
\hline
\end{tabular}
\end{table}

% -----------------------------------------------------------------------
\subsubsection{EV8 --- IOTA Transaction Submission}
\label{subsec:ev8-iota-transaction}

EV8 mencatat setiap pengiriman transaksi Move ke jaringan IOTA, baik yang
dilakukan oleh komponen \textit{client} (Patient, Hospital, Ministry) maupun
oleh PRE Server secara langsung menggunakan \texttt{proxy\_iota\_address}-nya
sendiri. Pencatatan \textit{event} ini diperlukan untuk mendeteksi modifikasi
transaksi selama pengiriman (T14).

\begin{table}[htbp]
\centering
\caption{Atribut Umum EV8 IOTA Transaction Submission}
\label{tab:ev8-umum}
\renewcommand{\arraystretch}{1.3}
\begin{tabular}{|p{3.5cm}|p{8.5cm}|}
\hline
\textbf{Atribut} & \textbf{Kemungkinan Nilai} \\
\hline
\texttt{source\_component} &
\texttt{HospitalTauriBackend} | \texttt{PatientTauriBackend} |
\texttt{MinistryTauriBackend} | \texttt{ProxyReencryption} \\
\hline
\texttt{actor\_id} &
IOTA \textit{address} penandatangan transaksi | \texttt{proxy\_iota\_address}
apabila PRE Server yang mengirim. \\
\hline
\texttt{target\_object} &
Nama fungsi Move yang dipanggil, misalnya \texttt{create\_access},
\texttt{signup}, \texttt{create\_medical\_record}. \\
\hline
\texttt{action\_type} &
\texttt{CREATE} \\
\hline
\texttt{outcome} &
\texttt{Success} | \texttt{Failure} \\
\hline
\end{tabular}
\end{table}

\begin{table}[htbp]
\centering
\caption{Atribut Tambahan EV8 IOTA Transaction Submission}
\label{tab:ev8-detail}
\renewcommand{\arraystretch}{1.3}
\begin{tabular}{|p{4cm}|p{7.5cm}|}
\hline
\textbf{Atribut} & \textbf{Kemungkinan Nilai dan Keterangan} \\
\hline
\texttt{transaction\_digest} &
\textit{Digest} transaksi IOTA yang dihasilkan setelah transaksi berhasil
dikirim dan dikonfirmasi. \\
\hline
\texttt{payload\_hash} &
Nilai \textit{hash} dari \textit{payload} transaksi sebelum dikirim ---
digunakan untuk mendeteksi modifikasi selama pengiriman. \\
\hline
\texttt{move\_function} &
Nama fungsi Move yang dipanggil, misalnya \texttt{create\_access},
\texttt{revoke\_access}, \texttt{create\_medical\_record}. \\
\hline
\texttt{move\_module} &
Nama modul Move tempat fungsi berada, misalnya \texttt{patient},
\texttt{proxy}, \texttt{admin}. \\
\hline
\texttt{network\_confirmation\_status} &
\texttt{confirmed} | \texttt{pending} | \texttt{failed} --- status
konfirmasi transaksi pada jaringan IOTA. \\
\hline
\texttt{gas\_used} &
Jumlah \textit{gas} yang dikonsumsi oleh transaksi. \\
\hline
\texttt{sponsor\_address} &
IOTA \textit{address} Gas Station yang mensponsori transaksi, apabila
transaksi menggunakan mekanisme sponsorship. \\
\hline
\end{tabular}
\end{table}

% -----------------------------------------------------------------------
\subsubsection{EV9 --- Gas Sponsorship Request}
\label{subsec:ev9-gas-sponsorship}

EV9 mencatat setiap permintaan \textit{sponsorship} transaksi kepada Gas
Station, baik dari komponen \textit{client} maupun dari PRE Server. Permintaan
ini mencakup dua tahap: reservasi \textit{gas object} (\texttt{/reserve\_gas})
dan eksekusi transaksi yang telah disponsori (\texttt{/execute\_tx}).
Pencatatan \textit{event} ini diperlukan untuk mendeteksi penyamaran sebagai
pemohon sah, penyangkalan atas permintaan \textit{sponsorship}, dan
pembanjiran \textit{gas object} (T15, T16, T17).

\begin{table}[htbp]
\centering
\caption{Atribut Umum EV9 Gas Sponsorship Request}
\label{tab:ev9-umum}
\renewcommand{\arraystretch}{1.3}
\begin{tabular}{|p{3.5cm}|p{8.5cm}|}
\hline
\textbf{Atribut} & \textbf{Kemungkinan Nilai} \\
\hline
\texttt{source\_component} &
\texttt{HospitalTauriBackend} | \texttt{PatientTauriBackend} |
\texttt{MinistryTauriBackend} | \texttt{ProxyReencryption} \\
\hline
\texttt{actor\_id} &
IOTA \textit{address} pemohon \textit{sponsorship} |
\texttt{proxy\_iota\_address} apabila PRE Server yang memohon. \\
\hline
\texttt{target\_object} &
\textit{Endpoint} Gas Station yang dituju: \texttt{/reserve\_gas} |
\texttt{/execute\_tx}. \\
\hline
\texttt{action\_type} &
\texttt{CREATE} | \texttt{EXECUTE} \\
\hline
\texttt{outcome} &
\texttt{Success} | \texttt{Failure} \\
\hline
\end{tabular}
\end{table}

\begin{table}[htbp]
\centering
\caption{Atribut Tambahan EV9 Gas Sponsorship Request}
\label{tab:ev9-detail}
\renewcommand{\arraystretch}{1.3}
\begin{tabular}{|p{4cm}|p{7.5cm}|}
\hline
\textbf{Atribut} & \textbf{Kemungkinan Nilai dan Keterangan} \\
\hline
\texttt{gas\_budget\_requested} &
Jumlah \textit{gas budget} yang diminta dalam permintaan \textit{sponsorship}. \\
\hline
\texttt{reserve\_duration\_secs} &
Durasi reservasi \textit{gas object} dalam detik sebelum kadaluarsa. \\
\hline
\texttt{reservation\_id} &
Pengenal unik reservasi \textit{gas} yang diterbitkan oleh Gas Station. \\
\hline
\texttt{sponsor\_address} &
IOTA \textit{address} Gas Station yang bertindak sebagai \textit{sponsor}. \\
\hline
\texttt{transaction\_digest} &
\textit{Digest} transaksi yang disponsori, tersedia setelah tahap eksekusi. \\
\hline
\texttt{gas\_coin\_object\_ids} &
Daftar ID objek \textit{gas coin} yang direservasi atau digunakan dalam
transaksi. \\
\hline
\end{tabular}
\end{table}

% -----------------------------------------------------------------------
\subsubsection{EV10 --- Redis Operation}
\label{subsec:ev10-redis-operation}

EV10 mencatat setiap operasi penyimpanan, pembacaan, atau penghapusan objek
sensitif pada Redis yang dilakukan oleh PRE Server. Objek yang dikelola
mencakup \textit{nonce} autentikasi (kunci Redis: \texttt{nonce:<iota\_address>})
dan material PRE berupa \textit{kfrag} serta \textit{access keys} (kunci Redis:
\texttt{keys:<hp\_addr>@<patient\_addr>}). Pencatatan \textit{event} ini
diperlukan untuk mendeteksi akses tidak sah atau modifikasi objek Redis (T18,
T19, T20).

\begin{table}[htbp]
\centering
\caption{Atribut Umum EV10 Redis Operation}
\label{tab:ev10-umum}
\renewcommand{\arraystretch}{1.3}
\begin{tabular}{|p{3.5cm}|p{8.5cm}|}
\hline
\textbf{Atribut} & \textbf{Kemungkinan Nilai} \\
\hline
\texttt{source\_component} &
\texttt{ProxyReencryption} \\
\hline
\texttt{actor\_id} &
\texttt{proxy\_iota\_address} PRE Server yang mengeksekusi operasi Redis. \\
\hline
\texttt{target\_object} &
Kunci Redis yang dioperasikan, misalnya
\texttt{nonce:<iota\_address>} atau
\texttt{keys:<hp\_addr>@<patient\_addr>}. \\
\hline
\texttt{action\_type} &
\texttt{CREATE} | \texttt{READ} | \texttt{DELETE} \\
\hline
\texttt{outcome} &
\texttt{Success} | \texttt{Failure} \\
\hline
\end{tabular}
\end{table}

\begin{table}[htbp]
\centering
\caption{Atribut Tambahan EV10 Redis Operation}
\label{tab:ev10-detail}
\renewcommand{\arraystretch}{1.3}
\begin{tabular}{|p{4cm}|p{7.5cm}|}
\hline
\textbf{Atribut} & \textbf{Kemungkinan Nilai dan Keterangan} \\
\hline
\texttt{redis\_key\_type} &
\texttt{nonce} | \texttt{access\_keys} --- kategori objek Redis yang
dioperasikan. \\
\hline
\texttt{operation\_type} &
\texttt{SET} | \texttt{GET} | \texttt{DEL} | \texttt{EXPIRE} ---
operasi Redis spesifik yang dieksekusi. \\
\hline
\texttt{ttl\_remaining} &
Sisa TTL entri Redis dalam detik pada saat operasi dilakukan;
\texttt{-1} apabila tidak ada TTL, \texttt{-2} apabila kunci tidak ditemukan. \\
\hline
\texttt{key\_pattern} &
Pola kunci Redis yang digunakan, misalnya \texttt{nonce:*} atau
\texttt{keys:*@*}, untuk membantu korelasi tanpa mengekspos nilai kunci
penuh. \\
\hline
\end{tabular}
\end{table}

% -----------------------------------------------------------------------
\subsubsection{EV11 --- IPFS Operation}
\label{subsec:ev11-ipfs-operation}

EV11 mencatat setiap operasi pengunggahan maupun pengunduhan objek pada IPFS
Cluster yang dilakukan oleh PRE Server. Objek yang dikelola adalah konten
klinis terenkripsi (AES-GCM) milik pasien. Pencatatan \textit{event} ini
diperlukan untuk mendeteksi pengunduhan data klinis oleh pihak tidak berwenang
maupun pembanjiran permintaan ke IPFS (T21, T22).

\begin{table}[htbp]
\centering
\caption{Atribut Umum EV11 IPFS Operation}
\label{tab:ev11-umum}
\renewcommand{\arraystretch}{1.3}
\begin{tabular}{|p{3.5cm}|p{8.5cm}|}
\hline
\textbf{Atribut} & \textbf{Kemungkinan Nilai} \\
\hline
\texttt{source\_component} &
\texttt{ProxyReencryption} \\
\hline
\texttt{actor\_id} &
\texttt{proxy\_iota\_address} PRE Server yang mengeksekusi operasi IPFS. \\
\hline
\texttt{target\_object} &
CID IPFS dari konten yang diunggah atau diunduh. \\
\hline
\texttt{action\_type} &
\texttt{CREATE} | \texttt{READ} \\
\hline
\texttt{outcome} &
\texttt{Success} | \texttt{Failure} \\
\hline
\end{tabular}
\end{table}

\begin{table}[htbp]
\centering
\caption{Atribut Tambahan EV11 IPFS Operation}
\label{tab:ev11-detail}
\renewcommand{\arraystretch}{1.3}
\begin{tabular}{|p{4cm}|p{7.5cm}|}
\hline
\textbf{Atribut} & \textbf{Kemungkinan Nilai dan Keterangan} \\
\hline
\texttt{cid} &
\textit{Content Identifier} IPFS dari objek yang dioperasikan. \\
\hline
\texttt{operation\_type} &
\texttt{upload} | \texttt{download} --- arah operasi terhadap IPFS. \\
\hline
\texttt{data\_size} &
Ukuran data dalam \textit{byte} yang diunggah atau diunduh. \\
\hline
\texttt{patient\_iota\_address} &
IOTA \textit{address} pasien pemilik data klinis yang dioperasikan. \\
\hline
\texttt{ipfs\_node\_url} &
URL \textit{node} IPFS Cluster yang digunakan dalam operasi. \\
\hline
\end{tabular}
\end{table}

% -----------------------------------------------------------------------
\subsubsection{EV12 --- IOTA Metadata Operation}
\label{subsec:ev12-iota-metadata}

EV12 mencatat setiap aktivitas pembacaan atau pembaruan objek \textit{on-chain}
pada jaringan IOTA, yang dapat dilakukan oleh seluruh komponen sistem. Objek
yang dicakup meliputi \texttt{AddressId}, \texttt{PatientIdAccount},
\texttt{HospitalPersonnelIdAccount}, \texttt{HospitalIdMetadata}, dan
objek-objek \textit{shared} lainnya. Pencatatan \textit{event} ini diperlukan
untuk mendeteksi potensi analisis pola transaksi dari metadata publik (T23).

\begin{table}[htbp]
\centering
\caption{Atribut Umum EV12 IOTA Metadata Operation}
\label{tab:ev12-umum}
\renewcommand{\arraystretch}{1.3}
\begin{tabular}{|p{3.5cm}|p{8.5cm}|}
\hline
\textbf{Atribut} & \textbf{Kemungkinan Nilai} \\
\hline
\texttt{source\_component} &
\texttt{HospitalTauriBackend} | \texttt{PatientTauriBackend} |
\texttt{MinistryTauriBackend} | \texttt{ProxyReencryption} \\
\hline
\texttt{actor\_id} &
IOTA \textit{address} pengguna | \texttt{proxy\_iota\_address} apabila
PRE Server. \\
\hline
\texttt{target\_object} &
IOTA Object ID dari objek \textit{on-chain} yang diakses. \\
\hline
\texttt{action\_type} &
\texttt{READ} | \texttt{UPDATE} \\
\hline
\texttt{outcome} &
\texttt{Success} | \texttt{Failure} \\
\hline
\end{tabular}
\end{table}

\begin{table}[htbp]
\centering
\caption{Atribut Tambahan EV12 IOTA Metadata Operation}
\label{tab:ev12-detail}
\renewcommand{\arraystretch}{1.3}
\begin{tabular}{|p{4cm}|p{7.5cm}|}
\hline
\textbf{Atribut} & \textbf{Kemungkinan Nilai dan Keterangan} \\
\hline
\texttt{object\_id} &
IOTA Object ID dari objek yang diakses atau diperbarui. \\
\hline
\texttt{object\_type} &
\texttt{AddressId} | \texttt{PatientIdAccount} |
\texttt{HospitalPersonnelIdAccount} | \texttt{HospitalIdMetadata} |
tipe objek lainnya. \\
\hline
\texttt{move\_function} &
Nama fungsi Move yang digunakan untuk mengakses objek, apabila akses
dilakukan melalui \texttt{dev\_inspect} atau transaksi. \\
\hline
\texttt{is\_mutable} &
\texttt{true} | \texttt{false} --- menyatakan apakah operasi bersifat
mutasi terhadap \textit{state on-chain} atau hanya pembacaan. \\
\hline
\texttt{transaction\_digest} &
\textit{Digest} transaksi IOTA apabila operasi bersifat mutasi;
\texttt{null} apabila hanya \texttt{dev\_inspect}. \\
\hline
\texttt{dev\_inspect\_used} &
\texttt{true} | \texttt{false} --- menyatakan apakah pembacaan dilakukan
menggunakan mekanisme \texttt{dev\_inspect} tanpa menerbitkan transaksi. \\
\hline
\end{tabular}
\end{table}

% -----------------------------------------------------------------------
\subsubsection{EV13 --- Capability Validation}
\label{subsec:ev13-capability-validation}

EV13 mencatat setiap proses validasi \textit{capability} yang dilakukan oleh
PRE Server sebelum permintaan sensitif diproses lebih lanjut. Objek yang
divalidasi meliputi \texttt{ProxyCapEnum} dan \texttt{GlobalAdminCap}
(\textit{on-chain}), serta JWT \textit{token} yang diterbitkan PRE Server untuk
\textit{hospital personnel}. Pencatatan \textit{event} ini diperlukan untuk
mendeteksi upaya eskalasi privilese melalui kelemahan validasi \textit{capability}
(T24, T25).

\begin{table}[htbp]
\centering
\caption{Atribut Umum EV13 Capability Validation}
\label{tab:ev13-umum}
\renewcommand{\arraystretch}{1.3}
\begin{tabular}{|p{3.5cm}|p{8.5cm}|}
\hline
\textbf{Atribut} & \textbf{Kemungkinan Nilai} \\
\hline
\texttt{source\_component} &
\texttt{ProxyReencryption} \\
\hline
\texttt{actor\_id} &
IOTA \textit{address} \textit{hospital personnel} yang divalidasi |
\texttt{proxy\_iota\_address} untuk validasi \texttt{ProxyCap}. \\
\hline
\texttt{target\_object} &
IOTA Object ID \textit{capability on-chain} | JWT \textit{token identifier}. \\
\hline
\texttt{action\_type} &
\texttt{VALIDATE} \\
\hline
\texttt{outcome} &
\texttt{Success} | \texttt{Denied} | \texttt{Failure} \\
\hline
\end{tabular}
\end{table}

\begin{table}[htbp]
\centering
\caption{Atribut Tambahan EV13 Capability Validation}
\label{tab:ev13-detail}
\renewcommand{\arraystretch}{1.3}
\begin{tabular}{|p{4cm}|p{7.5cm}|}
\hline
\textbf{Atribut} & \textbf{Kemungkinan Nilai dan Keterangan} \\
\hline
\texttt{capability\_id} &
IOTA Object ID \textit{capability on-chain} atau pengenal JWT yang
divalidasi. \\
\hline
\texttt{required\_scope} &
Cakupan akses yang dipersyaratkan oleh operasi yang diminta, misalnya
\texttt{read} atau \texttt{update}. \\
\hline
\texttt{actual\_scope} &
Cakupan akses aktual yang dimiliki oleh \textit{capability} yang
diperiksa. \\
\hline
\texttt{rejection\_reason} &
\texttt{expired} | \texttt{scope\_mismatch} | \texttt{not\_found} |
\texttt{invalid\_signature} | \texttt{null} apabila validasi berhasil. \\
\hline
\texttt{middleware\_layer} &
\texttt{jwt\_middleware} | \texttt{onchain\_check} --- lapisan validasi
yang melakukan pemeriksaan \textit{capability}. \\
\hline
\end{tabular}
\end{table}